%\subsection{Motivation}

% Motivação: contextualização do problema ou da necessidade que motivou o trabalho. Essa seção deve preparar o leitor para a <Pergunta da Pesquisa>

According to the World Health Organisation (WHO), there are at least 2.2 billion people with some visual impairment degree \cite{world2019world}. Among them, 43,3 million are classified as blind and 295 million have moderate or severe vision impairment. In order to be fully integrated into our society, they rely on assistive devices, such as canes, braille speakers, among others \cite{bourne2021trends}. 

Although a range of products has already been proposed, incorporating different features, they do not entirely fulfill their aim. Among the problems, of the solutions available in the market, are the lack of practicality and portability, invasive and requiring too much effort to learn \cite{lozano2009electrotactile}.

The difficulty of using or learn how to use a device could be avoided if concepts from Human Factors, or Ergonomics, were analysed during the product’s development, using appropriate methods. The early application of these methods and tests could be a gamechanger for the success of the product's user experience \cite{wolf2019towards}.

Motivated by the dissatisfaction of blind people with the currently available products, this paper starts from the hypothesis that a human-factors-centred design of assistive devices for blind and visually impaired people (BVIs) requires the involvement of BVIs in the design process in order to evaluate the product under design. The user has to test the product under development to provide feedback for the design team to improve the product.

In order to approach this problem, this work proposes using virtual reality (VR) as a tool for creating virtual environments, where proof of concepts or prototypes of assistive devices could be tested by BVIs. VR can be used to create specific, immersive and interactive situations that could help the user to learn and train \cite{farrell2018learning}, and the the developers to create more user-friendly products.

In a virtual environment, as long as the BVI is wearing a locating system, s/he can navigate the environment. Any information about the scenario, such as the position of objects and their distances to the user, is known and could be extracted from the virtual platform. As a consequence the designer can test different ways of translating this information into inputs before actually implementing a prototype of the assistive device, providing a flexible, safe and easy way to have it evaluated by different users.

%\section{Virtual reality in the design process}
\label{sec:vr_cabin}
The use of virtual reality for design purposes is not new. The cabin design process is often said to be complex because it involves several stakeholders, each with his/her own set of preferences and requirements. \cite{moerland2021application} proposed to anticipate the involvement of the final users based on co-design. In their proposal, the users can influence the product's development from the beginning. However, for the involvement to happen, a communication channel needed to be established, and it was done using virtual reality. The use case showed some benefits and disadvantages of using virtual reality. The virtual reality helped to bring the client closer to the design team, allowing them to draw quick sketches in brainstorming gatherings. It was associated with a steep learning curve for the designers. Among the disadvantages, it was considered a high-cost tool, and its use for a long time was associated with nausea.

%\section{Virtual reality for BVI users}
\label{sec:vr_without_vision}
Motivated by the popularization of virtual reality technology, \cite{siu2020virtual} developed a white cane to be used by BVI users in a virtual environment. Their purpose was to make virtual reality applications available for BVI users. In order to evaluate their proposal, the authors performed an experiment where the participants had to play a “scavenger hunt” using an HTC Vive system. Among the relevant findings of \cite{siu2020virtual} is that not all the participants reacted the same to a particular stimulus. The vibration of the cane was considered confusing by some participants, while others were familiar with it. Another interesting observation was that, similar to what happens in the real world, it was easier for the participants to navigate in larger areas than in tight spaces. Moreover, the authors observed that the participants focused their attention on the primary task, without freely exploring the environment, which might have impacted the low time to achieve the goal and the low number of obstacle hits. 

%\section{Augmented reality for BVI users}
\label{sec:ar_without_vision}
\cite{kirner2011using} raised two questions, "How can blind people learn 3D concepts aiming to be able to convert explored 3D environments into pictures?" and "How can we develop a spatial audio tutor with augmented reality technology to make easy the understanding of 3D concepts by blind people?" and used not using virtual reality technology but augmented reality to answer them. They developed a augmented reality application to be a tutor for BVI users. The application used allowed BVI users to play audio streams that were associeated with spatial positions. The users learned 3D concepts and also were able to  perceive, understand and produce embossed pictures representing real and imaginary 3D scenes. Also they were able to understand descriptions of 3D scenes described by non-BVI people. The authors believe that this application can be evolved to explain other concepts such as colors, transparency, shades, etc.

%\section{Information for BVI navigation}
\label{sec:bradley_dunlop}
Bradley and Dunlop published two works (\citeyear{bradley2002investigating,bradley2005experimental}) about how BVI navigates and how much it is similar or different to how a sighted person navigates. The first work of Bradley and Dunlop was published in \citeyear{bradley2002investigating} and discussed which type of information BVI uses to navigate in an environment and how it compares to sighted people. The second they compared the perceived workload of BVI participants and sighted participants when they navigate using user-tailored information created with the results of the previous experiments \cite{bradley2005experimental}. The results showed that BVI users reached landmarks significantly quicker when given the information made for that group, but still longer than sighted users. Also it showed that BVI participants systematically have a higher workload than sighted participants and that BVI users did have a higher workload when guided by orientations provided by sighted people, as well as the sighted participants did with orientations from BVI.

%\subsection{Mental Workload (MWL)}
\label{sec:mental_workload}
Mental workload is one of the main concepts studied in Human Factors \cite{stanton2004handbook}. The mental workload is similar to the physical workload but refers to the mental capacity necessary to perform a task. Each human being has a finite mental capacity. When the mental demand is higher than the operator's capacity, the person needs to adapt to finish the task, or the overall performance of the task is compromised. Otherwise, if the mental workload is too low, the operator may get bored and easily distracted and could also fail or not process the task's information. The mental workload is not a quantitative resource or something that one can directly measure, but several different techniques have been proposed in the literature to infer it and they can be: techniques based on task performance, techniques based on physiological measures and techniques based on subjective questionnaires.
        
%\begin{figure}[!htb]
    \centering

    \tikzstyle{arrow} = [rounded corners, line width = 1mm, ->]
    
    \resizebox{0.85\width}{!}{
    \begin{tikzpicture}[node distance=1cm]
        
        \node (mwl) {\includegraphics[width=.15\textwidth]{Fundamentação/Carga Mental/mwl.png}}
        node(t_mwl)[below of = mwl,yshift=-0.75cm] {Mental}
        node(t_mwl2)[below of = t_mwl,yshift=0.5cm] {workload};
        
        \node (demand) [above of=mwl, xshift=3cm, yshift=2.5cm] {\begin{tikzpicture}[node distance=1cm]
    \centering
    
    \node (worker) {\includegraphics[width=.15\textwidth]{Fundamentação/Carga Mental/work.png}};
    
    \node (energy1) [above of = worker, xshift = 0.5cm, yshift = 0.05cm]
     {\includegraphics[width=.060\textwidth]{Fundamentação/Carga Mental/bolt.png}} 
    node(energy2) [right of=energy1, xshift=-0.5cm, yshift = 0.3cm] {\includegraphics[width=.060\textwidth]{Fundamentação/Carga Mental/bolt.png}}
    node(energy3) [left of = energy1, xshift=0.5cm, yshift = 0.3cm] {\includegraphics[width=.060\textwidth]{Fundamentação/Carga Mental/bolt.png}};;
    
\end{tikzpicture}}
        node(t_demand)[below of = demand,yshift=-0.75cm] {Task}
        node(t_demand2)[below of = t_demand,yshift=0.5cm] {Demand};
        
        \node (capacity) [above of=mwl, xshift=-3cm, yshift=2.5cm] {\includegraphics[width=.15\textwidth]{Fundamentação/Carga Mental/full-battery.png}}
        node(t_capacity)[below of = capacity,yshift=-0.75cm] {Mental}
        node(t_capacity2)[below of = t_capacity,yshift=0.5cm] {Capacity};
        
        \node (tasks) [left of=mwl, xshift=-4cm, yshift=-6cm] {\includegraphics[width=.15\textwidth]{Fundamentação/Carga Mental/multitasking.png}}
        node(t_tasks)[below of = tasks,yshift = -0.75cm] {Primary and} 
        node(t_tasks2)[below of = t_tasks,yshift = 0.5cm] {secondary tasks};
        
        \node (physiological) [below of=mwl, yshift=-5cm,] {\includegraphics[width=.15\textwidth]{Fundamentação/Carga Mental/physiological.png}} 
        node(t_physiological) [below of = physiological, yshift = -0.75cm]{Physiological}
        node(t_physiological2) [below of = t_physiological, yshift = 0.5cm]{measurements};
        
        \node (subjective) [right of=mwl, xshift=4cm, yshift=-6cm] {\includegraphics[width=.15\textwidth]{Fundamentação/Carga Mental/subjective.png}}
        node(t_subjective) [below of = subjective, yshift = -0.75cm]{Subjective} 
        node(t_subjective2) [below of = t_subjective, yshift = 0.5cm]{measurements};
    
    
        \draw [arrow] (t_mwl2.south) to ++(0,-0.75) to +(-5,0) to (tasks.north);
        \draw [arrow] (t_mwl2.south) to (physiological.north);
        \draw [arrow] (t_mwl2.south) to ++(0,-0.75) to +(5,0) to (subjective.north);
        \draw [arrow, sharp corners] (capacity.east) to +(1.6,0) to (mwl.north);
        \draw [arrow, sharp corners] (demand.west) -- +(-1.35,0) to (mwl.north);
    
        
    \end{tikzpicture}
    }
    \caption{A overview of mental workload and the techniques to infer it.}
    \label{fig:mwl_overview}
\end{figure}    

%\subsubsection*{Techniques based on task performance}
%\label{subsec:task_performance}

%If the mental workload influences the task performance, it would be possible to infer it using the performance's variation of a task. Because there are cases where the user's mental capacity is too high for only one task, a common approach is to add a secondary task. In this case, the user is asked to maintain a good performance level and still try to execute both tasks. Both tasks should use the same skills \cite{stanton2004handbook, sanders1998human}.
%
%For example, an experiment to assess mental workload in a flight simulator may use two tasks. The primary task is to fly the aircraft while maintaining a performance level. The second is something simple, like mentally summing two random numbers that appear on the screen. If the numbers' sum is odd, the pilot should press the left arrow key on the keyboard, otherwise should press the right arrow key. If the pilot's performance in the secondary task is too low, it means that the mental demand from the first task is too high to pay attention to the second task \cite{mohanavelu2020cognitive}.
%
%%    \item Physiological measures;
%\subsubsection*{Techniques based on physiological measurements}
%%\label{subsec:physiological_measures}
%
%Many physiological measurements can be used to assess mental workload. The most common ones are heart and brain activity \cite{chakladar2020eeg, orlandi2018measuring}, skin conductance, eye movement and pupillary contraction \cite{stanton2004handbook, rodriguez2015pupillometry}. These measurements are considered an user-unbiased assessment technique \cite{fallahi2016effects}. It is recommended to evaluate them alongside another method, as they can be influenced by unknown variables and external factors. Some techniques use heart activities, obtained from an electrocardiogram (ECG) sensor, and electrodermal activity, obtained from a galvanic skin response (GSR) sensor, as physiological measurements to assess mental workload.
%
%The electrocardiogram (ECG) is a recording of the heart's electrical activity. From it, it is possible to determine the intervals between heartbeats and the corresponding frequency (heart rate, HR). Another common variable is the heart rate standard deviation (heart rate variability, HRV) \cite{cain2007review}. The heart activity is controlled by the sympathetic and parasympathetic nervous systems \cite{stanton2004handbook}. During a task, the heart activity changes with the mental demand of the task. The heart rate is expected to increase with the mental workload, while the heart rate variability is expected to decrease. These are consequences of two reactions in our system when in a mental demand situation \cite{stanton2004handbook}: a decrease in the parasympathetic nervous system activity and an increase in sympathetic nervous system activity. The ECG is a simple and non-invasive sensor used in many experiments to evaluate mental workload and other human factors' \cite{mohanavelu2020cognitive, mansikka2016fighter, zhang2014detection}.
%%(DEFINIR MELHOR)
%
%The skin electrodermal activity is affected by the person's sweating and the level of moisture in the environment. It can be used to reveal changes in our sympathetic system \cite{nourbakhsh2012using, shi2007galvanic}. It has been used in the literature as an assessment technique for stress and arousal \cite{nourbakhsh2012using, stanton2004handbook, shi2007galvanic}, the usability of human-computer systems \cite{shi2007galvanic} and also mental workload \cite{zhang2014detection, borghini2014measuring}.
%
%\subsubsection*{Techniques based on subjective questionnaires}
%%\label{subsec:subjective_measures}    
%
%The use of subjective questionnaires to assess mental workload has been extensively discussed in the literature \cite{sanders1998human, stanton2004handbook}. The questionnaires can be unidimensional or multidimensional. The unidimensional questionnaires are more straightforward by provide only a general workload score.     Examples of multidimensional questionnaires are the Subjective Workload Assessment Technique (SWAT) and the NASA Task Load Index (NASA-TLX). SWAT decomposes the mental in three dimensions: time load, mental effort load, and psychological stress. NASA-TLX, a questionnaire created by \cite{hart1988development}, uses six dimensions, as described in Table \ref{tab:nasa_dimensions}. These questionnaires were proposed to evaluate only one task/activity. If the user has performed two tasks (a primary and secondary task), he/she should be oriented to answer about the primary task, not a combination of both \cite{sanders1998human}.
%
%\begin{table}[htb]
%    \centering
%    \caption{NASA-TLX dimensions and the description of each dimension. \cite{stanton2004handbook}.}
%    \label{tab:nasa_dimensions}
%    \begin{tabular}{|l|l|}
%
%        \hline
%        \textbf{Dimension}   & \textbf{Explanation}                                                                                                                                                       \\ \hline
%        Mental demand (MD)   & \begin{tabular}[c]{@{}l@{}}The mental and perceptive activity\\ demanded by the task (chose, decide,\\ think, calculate, search, etc.).\end{tabular}                       \\ \hline
%        Physical demand (PD) & \begin{tabular}[c]{@{}l@{}}The physical activity demanded by\\ the task (pull, lift, spin, drag, etc.).\end{tabular}                                                       \\ \hline
%        Temporal demand (TD) & \begin{tabular}[c]{@{}l@{}}The time pressure felt by the user.\\ A rating the leverages the time \\ available and the time necessary to\\ completed the task.\end{tabular} \\ \hline
%        Performance (PE)     & \begin{tabular}[c]{@{}l@{}}The user's satisfaction with its \\ performance or result the task.\end{tabular}                                                                \\ \hline
%        Effort (EF)          & \begin{tabular}[c]{@{}l@{}}A rating of the effort necessary \\ to achieve that performance felt by\\ the user.\end{tabular}                                                \\ \hline
%        Frustration (FR)     & \begin{tabular}[c]{@{}l@{}}A rating of stress, annoy or irritation\\ felt by the user throughout the task.\end{tabular}                                                    \\ \hline
%    \end{tabular}
%\end{table}
%
%Finally, it is essential to highlight that, to have a comprehensive evaluation of the mental workload, not to choose only one technique, but to combine techniques from the three classes (task performance, physiological measurements and subjective questionnaires). Mental workload is multidimensional and can reflect partially or differently in each technique \cite{sanders1998human}.

%\subsection{Situation Awareness (SA)}
\label{sec:situation_awareness}
%Situation awareness can be defined as “the perception of the elements within a volume of time and space (Level 1), the comprehension of their meaning (Level 2), and the projection of their status in the near future (Level 3)” as illustrated in Figure \ref{fig:sa_overview}. One example is when an air traffic controller looks at a radar display (Level 1). He/she seeks to understand the aircraft's position and speed (Level 2) and then predict its position in the near future, 5, 10 or 15 minutes after (Level 3) \cite{sanders1998human}. Similarly, when a pilot reads the cockpit panel (Level 1) and understands their data (Level 2) then he/she can predict the next reading of that same instrument or some other status of the aircraft after a couple of minutes (Level 3).

%\begin{figure}[!htb]
    \centering
    \tikzstyle{arrow} = [rounded corners, line width = 1mm, ->]

    \resizebox{0.85\textwidth}{!}{
    \begin{tikzpicture}[node distance=3.3cm]
        \centering
        
        \node (info) [fill=white] {\includegraphics[width = 0.20\textwidth]{Fundamentação/Percepção situacional/information.png}}
        node(t_info)[below of = info,yshift=1.3cm] {\Large Information};
        
        \node (arrow) [right of=info, xshift=7.0cm, yshift = -0.5cm] {\begin{tikzpicture}[node distance=5cm, line width = 2.5mm]
    \centering

    \node (seta1) [] {};
    \node (seta2) [right of = seta1, xshift = 5.0cm] {};
    \node (seta3) [above of = seta2, yshift = -2.5cm] {};
    \node (seta5) [below of = seta1, yshift = -1.5cm,] {};
    \node (seta6) [right of = seta5, xshift = 5.0cm] {};
    \node (seta7) [below of = seta6, yshift = 2.5cm] {};
    \node (seta4) [right of = seta1, below of = seta1, yshift = 1.75cm, xshift = 10.0cm] {};

    \draw (seta1.center) -- (seta2.center);
    \draw (seta2.center) -- (seta3.center);
    \draw (seta3.center) -- (seta4.center);
    \draw (seta4.center) -- (seta7.center);
    \draw (seta7.center) -- (seta6.center);
    \draw (seta6.center) -- (seta5.center);
    \draw (seta5.center) -- (seta1.center);
    

\end{tikzpicture}
};
        
        \node (perception) [right of=info, xshift=2.5cm, draw, line width=2mm] {\begin{tikzpicture}[node distance=5cm]
    \centering
    
    \node (comprehension) {\includegraphics[width=3.0cm]{Fundamentação/Percepção situacional/eye.png}};

    \node(info) [] {\includegraphics[width=1.0cm]{Fundamentação/Percepção situacional/information}};

\end{tikzpicture}}
        node(t_perception)[below of = perception, yshift=0.9cm] {\Large 1st Level}
        node(t_perception2)[below of = t_perception, yshift=2.5cm] {\Large Perception};

        \node (comprehension) [right of=perception, xshift=0.25cm, draw, line width=2mm] {\begin{tikzpicture}
    \centering
    
    \node (comprehension) {\includegraphics[width=3.0cm]{Fundamentação/Percepção situacional/idea.png}};
    
    \node(info) [yshift = 0.5cm] {\includegraphics[width=0.75cm]{Fundamentação/Percepção situacional/information}};

\end{tikzpicture}}
        node(t_comprehension)[below of = comprehension, yshift=0.9cm] {\Large 2nd Level}
        node(t_comprehension2)[below of = t_comprehension, yshift=2.5cm] {\Large Comprehension};

        \node (projection) [right of=perception, xshift=3.75cm, draw, line width=2mm] {\begin{tikzpicture}[node distance=5cm]
    \centering
    
    \node (comprehension) {\includegraphics[width=3.0cm]{Fundamentação/Percepção situacional/future.png}};

    \node(info) {\includegraphics[width=1.5cm]{Fundamentação/Percepção situacional/information}};

\end{tikzpicture}
}
        node(t_projection)[below of = projection, yshift=0.9cm] {\Large 3rd Level}
        node(t_projection2)[below of = t_projection, yshift=2.5cm] {\Large Projection};

        \node (decision) [right of=projection, xshift=5.0cm] {\includegraphics[width = 0.30\textwidth]{Fundamentação/Percepção situacional/decision.png}}
        node(t_decision)[below of = decision, yshift=0.5cm] {\Large Decision}
        node(t_decision2)[below of = t_decision, yshift=2.5cm] {\Large making};

    \end{tikzpicture}
    }
    \caption{An overview of situation awareness and the SAGAT.} %and the methods to infer it
    \label{fig:sa_overview}
\end{figure}

The term “situation awareness” was first proposed for the Aeronautics domain and today is considered a key factor for designing complex and dynamic systems from other domains, such as automotive, medical and nuclear \cite{endsley1995measurement}. It can be defined as “the perception of the elements within a volume of time and space (Level 1), the comprehension of their meaning (Level 2), and the projection of their status in the near future (Level 3)” \cite{sanders1998human}. It is an essencial factor to make sure that the user will be capable to make important decisions correctly and achieve high-performance \cite{endsley1988design, endsley2018automation}. As it is for the mental workload, situation awareness is not a quantitative subject. The most common way to measure it is using subjective techniques, among which one of the most famous is the Situation Awareness Global Assessment Technique (SAGAT). It was proposed by \cite{endsley1988design} and is based on how the information is processed inside the user’s mind.

%The test application is made by freezing the operator activity, usually made in a simulation environment, and then asking the user some questions that were previously defined based on the user's activity. These questions should be as similar as possible to how the person thinks when reasoning about the situation to avoid extra effort in understanding it \cite{stanton2004handbook}.  Although freezing the activity may sound troublesome, empirical work has shown that it does not interfere with the user performance and the user memory can withstand a break for as long as 5 to 6 min \cite{endsley1988design}.


%\section{Co-Design}
\label{sec:co_design}
Co-design, or collaborative design, refers to a design process in which individuals of the design team have different backgrounds or bring different experiences, which can be essencial for the product under design. It is based on good communication and information sharing among the team \cite{chiu2002organizational}.

%\cite{kleinsmann2006understanding} provides the following definition:.
%
%\begin{quote}
%    \textit{"Collaborative design is the process in which actors from different disciplines share their knowledge about both the design process and the design content. They do that to create a shared understanding of both aspects, to be able to integrate and explore their knowledge and to achieve the larger common objective: the new product to be designed."} \cite{kleinsmann2006understanding}.
%\end{quote}
%
%This definition emphasizes two critical aspects of co-design: knowledge sharing and integration. According to \cite{kleinsmann2006understanding} knowledge is the data after the receiver's understanding or translating process, in a state that is possible to record or register, so that the person can remember and use it later. During the collaborative design, ideas, facts or concepts are exchanged between the actors. This exchange is a fundamental part of the co-design process since it is responsible for the growth of each individual's knowledge. Once the knowledge is shared among the actors, they can use it when performing their tasks, resulting in knowledge integration \cite{kleinsmann2006understanding}.

This paper's main goal is the use of virtual reality as a tool for evaluating proofs of concept of assistive devices for blind and visually impaired people from a human-factors perspective. The purpose is to provide a flexible and easily configured way of testing different concepts of assistive devices in order to support an agile and user-centered development.

This goal is related to the following research questions, which are investigated in this work:

\begin{itemize}
   \item Is it possible to evaluate and compare concepts of assistive devices from a human factors perspective in a virtual environment? What are the main limitations of the use of a virtual reality environment? \label{itm:obj_first}
   \item Do non-BVI users, when deprived of their vision, similarly evaluate assistive devices as BVI users? \label{itm:obj_second}
\end{itemize}

% Delimitação da pesquisa: é o recorte do seu trabalho

The concepts of assistive devices presented as part of this work are used only as examples for investigating the research questions presented. The challenges related to their full development up to high Technology Readiness Levels (TRLs), as well as their feasibility as commercial products, are out of the scope of this work.

% Estrutura do texto

\subsection*{Structure of the text}

The next Section of this paper are organized as follows.

Section \ref{ch:metodologia} details the proposal of this paper describing how virtual reality could be used to integrate BVI users into the design process of assistive design. It illustrates the proposed method by applying it to evaluate three different assistive devices (audio guide, virtual cane and haptic belt), as well as their mixed-use, in the environment of a hospital reception. 

Section \ref{ch:resultados} describes the experiment designed to evaluate the paper's proposal and analyses the results in order to investigate the research questions of Section \ref{sec:objetivos}

Finally, section \ref{ch:conclusao} summarizes the main conclusions of this work and discusses future work.
